\begin{trapbox}

	\begin{description}[style=nextline, leftmargin=2.8em, labelsep=0.5em, itemsep=0.35em]
		\item[\textbf{\textcolor{CTrap}{易错点一（公比为零）}}] 在使用公式 $a_n = a_m q^{n-m}$ 时，忽略公比 $q$ 可能为零。
		\item[\textbf{\textcolor{CTrap}{正确理解（非零前提）：}}] 等比数列要求 $q \neq 0$，否则数列退化或运算无意义。
		\item[\textbf{\textcolor{CTrap}{错因分析（忽略定义）：}}] 未注意等比数列的定义条件，直接套公式。
	\end{description}

	\medskip
	\begin{description}[style=nextline, leftmargin=2.8em, labelsep=0.5em, itemsep=0.35em]
		\item[\textbf{\textcolor{CTrap}{易错点二（指数方向）}}] 误写指数为 $m-n$ 而不是 $n-m$。
		\item[\textbf{\textcolor{CTrap}{正确理解（下标差规则）：}}] 指数是 $n-m$，不要写成 $m-n$。当 $n<m$ 时指数为负，相当于除以 $q$ 的正次幂。
		\item[\textbf{\textcolor{CTrap}{错因分析（顺序混淆）：}}] 未理清“从 $a_m$ 到 $a_n$ 需要乘几次 $q$”。
	\end{description}

	\medskip
	\begin{description}[style=nextline, leftmargin=2.8em, labelsep=0.5em, itemsep=0.35em]
		\item[\textbf{\textcolor{CTrap}{易错点三（负公比符号）}}] 计算负数公比的幂时，忽视奇偶性导致正负号错误。
		\item[\textbf{\textcolor{CTrap}{正确理解（奇负偶正）：}}] $(-2)^3 = -8$，$(-2)^4=16$，不能简单地取绝对值运算。
		\item[\textbf{\textcolor{CTrap}{错因分析（指数运算习惯）：}}] 常用正数公比，忽略负数时需逐层确定符号。
	\end{description}

\end{trapbox}
